\section{System Implications and Research Outlook}
\label{sec:hardwareconclusion}
The prediction studies select reproduction programs; reference gates decide compatibility; cost studies bound optimization opportunity (Table~\ref{tab:decisions}). B570's successful simplification removes draft/proposal dependencies from its frozen reproduction task. RTX's specified capacity-level simplification fails, while proposal-conditioned replay reproduces the complete path. Neither result establishes a minimal state representation or a portable cause.
\begin{table*}[t]
\caption{Engineering uses of the retained results. Reproduction concerns frozen outputs; optimization budgets are conditional and no proposed improvement is measured.}
\label{tab:decisions}
\centering
\begin{tabular}{@{}>{\raggedright\arraybackslash}p{.62in}>{\raggedright\arraybackslash}p{1.58in}>{\raggedright\arraybackslash}p{2.22in}>{\raggedright\arraybackslash}p{1.68in}@{}}
\toprule
Study & Information beyond a basic gate & Reproduction task / optimization decision & Scope of that use \\
\midrule
B570 & Full non-reference suffixes predicted without proposals & Reproduce the tested capped native outputs with target-only UB1; remove draft loading, proposals and native speculative outer loop & Frozen stack and prompts; no measured debugging saving or causal irrelevance \\
RTX A & Candidates with equal totals fail on different tasks; replay succeeds & Reject this UB2 reproducer; retain the sealed tape, locally computed target decisions and KV reconciliation & Shared numerical routines; no minimality or validation of a correction \\
RTX B & Serialization qualifies but loses on every prompt & Exclude measured serialization as a latency optimization & Instrumented post-prefill cost; no qualified batched gain \\
358H & Overhead-only prompt-GM ceilings remain below one & Parity needs a changed premise; even with other cost zero, uniform target-work cuts must reach \DiagCPUReduction\%/\DiagNPUReduction\% & Baseline/progress fixed; necessary optimistic budgets, not achieved savings \\
\ifruntime M4 & About \DiagMOpportunity\% two-fixed-mode opportunity; Guard adds 2.71\% GM latency & Prefer fixed speculation on this aggregate measured workload; assess opportunity before control & Hindsight fixed modes; no bound on arbitrary within-request switching \\
\fi
\bottomrule
\end{tabular}
\end{table*}


\paragraph{Use a tested reproducer for diagnosis.}
On B570, a target-only harness can generate the tested native expected outputs, including non-reference suffixes. On RTX, the successful harness computes its own target decisions while consuming a sealed proposal tape. The complete arrays give expected outputs for regression checks of the frozen paths. Their agreement supplies no evidence that a future correction or changed numerical stack will behave identically; such a use needs agreement back to native execution. The present study demonstrates reproducer selection, not a completed bug fix or measured debugging savings.

\paragraph{Change the cost term that limits the opportunity.}
RTX serialization qualifies but loses within its instrumented boundary. On 358H, removing draft/residual cannot reach parity at fixed target work and progress. Under the optimistic premise that all other cost vanishes, a uniform target-work reduction of at least \DiagCPUReduction\% CPU or \DiagNPUReduction\% NPU is needed for the original prompt-GM ceiling to reach one. A design that changes target execution or committed progress must qualify that new path before claiming a gain. Isolated NPU speed cannot establish it.

\ifruntime
\paragraph{Measure opportunity before adding control.}
Hindsight choice between the two measured fixed modes lowers GM latency from Always by about \DiagMOpportunity\%; Guard instead adds 2.71\%. This small fixed-mode opportunity helps explain why fallback activity alone is weak evidence of useful control. It does not bound policies that change the execution or switch within a request.
\else
The full manuscript separately assesses M4 control and its limited opportunity between two measured fixed modes.
\fi

A useful next diagnostic study would test whether a simpler reproducer helps isolate a discrepancy or validate a correction, then confirm the result on native execution. Stronger causal attribution instead needs a discriminating state/trace-matched contrast and controlled continuation. These are additional research questions; full-array agreement alone cannot answer them. The conclusions here remain the tested reproduction choices and conditional cost budgets.
